\section{Descrição do Problema}\label{sec:descricao_do_problema}

A identificação visual de sítios ligantes em proteínas é uma etapa imprescindível no desenvolvimento de fármacos. Entretanto, os softwares de visualização molecular convencionais apresentam uma limitação estrutural: por operarem, majoritariamente, em dispositivos de apresentação planos, não oferecem ao usuário a sensação de profundidade necessária para compreender a geometria tridimensional do sítio de ligação, prejudicando a formação de uma noção espacial precisa da molécula.

Embora ferramentas com suporte à realidade virtual representem uma alternativa promissora a esse problema, as opções disponíveis atualmente são escassas e, em sua maioria, dependentes de hardware dedicado de alto desempenho gráfico — o que restringe o acesso a usuários com computadores de menor capacidade computacional. Ademais, essas soluções geralmente exigem instalação local, limitando ainda mais sua adoção em contextos acadêmicos e científicos com infraestrutura reduzida.

Diante desse cenário, identifica-se a necessidade de uma solução de visualização molecular em realidade virtual que seja leve, acessível e executável diretamente pelo navegador, sem dependência de hardware especializado, democratizando o acesso a essa forma de visualização.